% 编译: xelatex model.tex   (中文需要 XeLaTeX + ctex)
\documentclass[11pt,a4paper]{article}
\usepackage{ctex}
\usepackage{amsmath,amssymb}
\usepackage{booktabs}
\usepackage{array}
\usepackage{geometry}
\geometry{margin=2.4cm}
\usepackage{enumitem}
\usepackage[colorlinks=true,linkcolor=blue,citecolor=blue]{hyperref}

\newcommand{\BW}{\mathrm{BW}}
\newcommand{\MAC}{\mathrm{MAC}}
\newcommand{\LB}{\mathrm{LB}}
\newcommand{\TILEM}{\mathrm{TILE\_M}}
\newcommand{\TILEN}{\mathrm{TILE\_N}}
\newcommand{\ACTHALF}{\mathrm{ACT\_HALF}}
\newcommand{\code}[1]{\texttt{#1}}

\title{MoE 算子 cost model: 理论建模}
\author{}
\date{}

\begin{document}
\maketitle

\begin{abstract}
本文给出一个昇腾 NPU 上 MoE 融合算子 cost model 的理论建模。

\textbf{输入}分四组事实: \emph{形状}给出模型维度与负载, 后者是一个路由计数张量
$C[\text{dst}][\text{expert}][\text{src}]$ 而不只是 token 总数; \emph{硬件}给出核数与各级
带宽、容量的规格与标定值; \emph{编译点}给出 tile 几何、L1 组织、量化格式与通信模式这些
决定二进制的轴; \emph{编排}给出波如何推进、每个 stage 落在哪个执行角色、tile 以多大粒度
成为一次工作、中间结果落片上还是落显存、tile 到核的绑定在编译期定死还是推迟到派发时刻。

\textbf{原理}是不拟合端到端公式, 而是构造并求解一个资源受限调度实例, 并把\emph{成本}与
\emph{约束}分开表达: 成本在降解期逐 tile 产生 —— 一份工作的时长取其计算量与搬运量的较大
者, 是局部的 roofline 而非对整个算子套一次; 约束写进事件的五个字段 —— 独占资源、带延迟的
前驱、跨事件持有的计数额度、共位锚点与时长; 时间则完全在调度期产生, 由表调度在这些约束下
摆放事件得到。因此``算力绑定还是带宽绑定''是涌现的结论而非假设的前提。算子在给定形状、
硬件、编译点与编排下被降解成一张带标注的事件图, 执行时间是该图在一套调度规则下的
makespan。编排因此是模型的\textbf{输入}而非求解目标 —— 给定一种编排, \textbf{输出}是该
编排下的执行时间与一份可追问的分解: 逐 stage 的忙碌、跨度与三类等待 (依赖、资源、额度),
关键路径及其上每一跳的阻塞成因, 三条只由形状决定的物理下界与当前被哪一条绑定, 以及每个
角色池上的空闲按``无活可做''与``有活可做却空着''的切分。模型的主产品是同一份声明下改动
单个编排或 tiling 旋钮后的时间\emph{差值}: 两次运行共享全部未标定常数, 系统误差同向抵消,
并各自通过同一组护栏, 故差值可比而非模拟器的假象。

全文的数值实例取 $h=5120$、中间维 $I=4608$、每卡 72 token、$\mathrm{topk}=6$、8 卡每卡
3 专家、28 个 Cube 核。
\end{abstract}

\tableofcontents

%======================================================================
\section{建模输入}
\label{sec:input}

一次运行是一个实例
\begin{equation}
I = I(\underbrace{\text{形状}}_{\text{算什么}},\ \underbrace{\text{硬件}}_{\text{什么机器}},\
      \underbrace{\text{编译点}}_{\text{什么二进制}},\ \underbrace{\text{编排}}_{\text{怎么组织}}),
\label{eq:instance}
\end{equation}
四样缺一样, ``跑多久''就没有唯一答案。四组输入的内容如下。

\begin{center}
\begin{tabular}{@{}llp{7.1cm}@{}}
\toprule
组 & 性质 & 内容 \\
\midrule
形状 & 任务事实
 & 模型维度 ($h$, 中间维 $I$); 负载为路由计数张量
   $C[\text{dst}][\text{expert}][\text{src}]$, 不只是 token 总数 ——
   专家间的不均衡由它直接表达。守恒性 (每源卡发出
   $\text{tokens}\times\text{topk}$ 行) 是硬约束 \\
硬件 & 规格与标定
 & 核数 $P$; 各级带宽 (显存$\to$L1、片上向量、跨卡读、跨卡写) 与容量
   (L1 / UB / L0C); Cube 速率。每个常数带出处标签, 分为规格、实测、
   实现取值与假设四类 \\
编译点 & 二进制事实
 & tile 几何 ($\TILEM$, $\TILEN$, L1 的 $K$ 分块)、L1 缓冲块数、权重布局与
   是否交织、量化模式与数据类型、通信模式。十余个轴取其规范哈希为
   \emph{编译指纹} —— 内核自身的 tiling key 只编码其中一部分, 故同一个 key
   可对应多个二进制 \\
编排 & \textbf{可扫的选择}
 & 波如何推进 (波宽与波序列); 每个 stage 落哪个执行角色; tile 以多大粒度
   成为一次工作; stage 之间那条边的就绪粒度与中间结果落点 (片上/显存);
   tile 到核的绑定在降解期定死还是推迟到派发时刻; 就绪队列的排序规则 \\
\bottomrule
\end{tabular}
\end{center}

\paragraph{编排是输入, 不是求解目标.} 这是全文的立足点。许多 cost model 把``最优编排''
当作待求解的对象; 本模型不求解它 —— 给定一种编排, 模型给出它的代价, 比较两种编排就跑
两次。直接后果是模型的主产品为\emph{差值} $\Delta\widehat{T}$ 而非单次的绝对时间
(第 \ref{sec:use} 节)。

\paragraph{输入里刻意不含什么.}
形状与拓扑\textbf{不进编译指纹}: 它们每次运行都变, 混入后指纹就失去``同一个二进制''
的含义, 二者分属形状域与运行拓扑。输入里也没有任何拟合系数: 每个常数都是机制量
(带宽、速率、容量、容量深度), 且一个参数必须具有可建模的物理后果才被接受 ——
否则对它的扫描只会得到``改它没用''的假结论。

\paragraph{输入进入模型的位置.} 实现上分六层: 输入事实 $\to$ 适配器 (编排) $\to$ 建图器
(逐 tile 展开) $\to$ 事件图 $\to$ 调度器 $\to$ 分析。形状与编译点在第一、三层被消耗,
编排在第二层被声明并在第三、五层生效, 硬件常数在第三层 (成本) 与第六层 (下界) 各用一次。
\textbf{成本在建图期产生, 时间在调度期产生}, 这是与解析型 cost model 的根本区别; 混淆
两者是这类模型最常见的错误 —— 例如把贪心调度的结果当成硬件行为。

%======================================================================
\section{形状算术: workload $\to$ tile $\to$ wave}
\label{sec:shape}

\subsection{路由}
每个源卡发出 $\text{tokens}\cdot\text{topk}$ 行。均匀路由下每个专家收到
\begin{equation}
\text{rows}_e = \frac{\text{tokens}\cdot\text{topk}\cdot W}{E_{\text{total}}}
             = \frac{72\times 6\times 8}{24} = 144 .
\label{eq:rows}
\end{equation}
非均匀路由由路由计数张量 $C[\text{dst}][\text{expert}][\text{src}]$ 精确给出, 守恒性
(每源卡发出 $\text{tokens}\times\text{topk}$ 行) 是硬约束, 违反即报错。

\subsection{tile 网格}
行方向按 $\TILEM$ 切成 m-group:
\begin{equation}
n_m = \left\lceil \text{rows}_e / \TILEM \right\rceil = \lceil 144/256\rceil = 1 .
\end{equation}
列方向 GMM1 与 GMM2 的 tile 数\textbf{不同}:
\begin{align}
n^{(1)}_{\text{tile}} &= \left\lceil \frac{\lceil \mathrm{hidden\_dim}/\ACTHALF \rceil}{\TILEN} \right\rceil
   = \lceil 4608/256 \rceil = 18, \label{eq:n1}\\
n^{(2)}_{\text{tile}} &= \left\lceil \frac{h}{\TILEN} \right\rceil
   = \lceil 5120/256 \rceil = 20 . \label{eq:n2}
\end{align}
\eqref{eq:n1} 的调度宽度取 $\mathrm{hidden\_dim}/\ACTHALF$ 而非 $\mathrm{hidden\_dim}$:
非交织布局下一个 tile 要跑两遍权重块 (SwiGLU 的 gate 与 up)。GMM2 的收缩维
\begin{equation}
K_2 = \mathrm{hidden\_dim}/\ACTHALF = 4608 \quad (=\text{中间维} I).
\end{equation}

\subsection{波宽}
一波覆盖多少个 m-group:
\begin{equation}
\mathrm{mgw} = \max\left(
  \left\lceil \frac{P\,p_1}{\lceil \mathrm{hidden\_dim}/\TILEN \rceil} \right\rceil,\
  \left\lceil \frac{P\,p_2}{\lceil h/\TILEN \rceil} \right\rceil \right),
\label{eq:mgw}
\end{equation}
$P$ 为 Cube 核数, $p_1,p_2$ 是``每核至少领到几个逻辑 tile''的策略值。本例
$\mathrm{mgw} = \max(\lceil 56/36\rceil, \lceil 28/20\rceil) = 2$, 共 2 波。
\eqref{eq:mgw} 的含义是\textbf{用波宽保证并行度}: 波太窄则核领不满, 波太宽则下游就绪太晚。

%======================================================================
\section{单事件代价}
\label{sec:cost}

统一形式是逐 tile 的 roofline:
\begin{equation}
d(e) = \max\left(\frac{\MAC(e)}{R_{\text{cube}}},\ \frac{\mathrm{bytes}(e)}{\BW}\right).
\label{eq:roofline}
\end{equation}
取 $\max$ 而非求和, 是``搬运与计算完全重叠''的理想化, 方向偏乐观 —— 因此第
\ref{sec:bounds} 节的下界检查才有意义 (偏小的模型撞下界说明输入错了)。注意
\eqref{eq:roofline} 是\textbf{局部}的: 不对整个算子套一次 roofline, 全局的``算力绑定
还是带宽绑定''是从调度中涌现的结果。

\subsection{dispatch: 行级软流水}
\begin{equation}
T_{\text{seg}} = \max\!\big(\lambda,\ \min(\text{rows}, d_{\text{buf}})\, t_{\text{row}}\big)
               + \max(0,\ \text{rows}-d_{\text{buf}})\, t_{\text{row}} ,
\label{eq:dispatch}
\end{equation}
$d_{\text{buf}}$ 为软流水槽数 (内核 \code{dispatchBufferCount}, 取 6),
$t_{\text{row}} = \text{每行字节}/\BW$, $\lambda$ 为一次往返延迟 (本卡/跨卡不同)。
$\text{rows}\le d_{\text{buf}}$ 时段时长是一个\textbf{平台}, 不随行数增长: 前
$d_{\text{buf}}$ 行的取数背靠背发出, 只有在途数达到槽数才等待腾槽。朴素式
$\lambda + \text{rows}\cdot t_{\text{row}}$ 对实测高估 32\%。

\subsection{GMM1}
\begin{align}
L_1 &= \underbrace{\frac{m K}{\BW}}_{\text{A 流 (激活)}}
     + \underbrace{\frac{w_b\, K\, \text{cols}}{\BW_b}\, f}_{\text{B 流 (权重)}},
\qquad
C_1 = \frac{2\, m\, \text{cols}\, K}{R_{\text{cube}}}, \label{eq:gmm1-phase}\\
T_{\text{gmm1}} &= \begin{cases}
  \max(L_1, C_1), & b \ge 2 \quad (\text{双缓冲}),\\
  L_1 + C_1 + n_{\text{chunk}}\, t_{\text{restart}}, & b = 1 .
\end{cases}
\label{eq:gmm1}
\end{align}
$w_b$ 为权重块数 ($=\ACTHALF$, 交织布局下为 1), $f$ 为 B 流复用比例,
$n_{\text{chunk}} = \lceil K/\mathrm{l1\_tile\_k}\rceil$。系数 2 来自 SwiGLU 双投影。

\paragraph{两股流相加而非取 $\max$.} 这一点由对照实验判定, 源码读不出来。固定并发核数
只变 $m$ 时, $\max$ 口径预测斜率为 0 (因 B 流恒大于 A 流), 而实测两对比较分别相差
36\% 与 86\%; 另有一对固定 $m$ 只变每专家 m-group 数的比较, 说明 $\max$ 在大 batch 上
的那点支持来自 B 流复用 (L2 命中), 不是重叠。

\subsection{activation}
\begin{equation}
T_{\text{act}} = T_{\text{startup}} + \frac{m\,\text{cols}}{V}\cdot\frac{\beta_{\text{vec}}}{\BW_{\mathrm{UB}}},
\qquad V = 64,\quad \beta_{\text{vec}} = 722\ \mathrm{B} .
\label{eq:act}
\end{equation}
$\beta_{\text{vec}}$ 由源码逐句计数得到: bf16 中间缓冲被整体流三遍 (SwiGLU 写一遍,
求最大指数读一遍, 量化再读一遍)。两点 $m$ 扫独立验证: 实测斜率 $0.007760\ \mu s$/向量,
而 $\beta_{\text{vec}}/\BW_{\mathrm{UB}} = 0.007763$ (偏差 $+0.05\%$)。

\subsection{GMM2}
\begin{align}
L_2 &= \underbrace{\frac{m K_2}{\BW}\,\mathbb{1}[\text{物化}]}_{\text{A 流}}
     + \underbrace{\frac{K_2\,\text{cols}}{\BW_b}}_{\text{B 流}},
\qquad C_2 = \frac{m\,\text{cols}\,K_2}{R_{\text{cube}}}, \label{eq:gmm2-phase}\\
T_{\text{gmm2}} &= \begin{cases}\max(L_2,C_2), & b\ge 2,\\ L_2+C_2, & b=1.\end{cases}
\label{eq:gmm2}
\end{align}
与 \eqref{eq:gmm1} 的三处差异各有不同性质:

\begin{center}
\begin{tabular}{@{}lll@{}}
\toprule
差异 & 根源 & 性质 \\
\midrule
无系数 2 & 单投影 & 算法事实 \\
无 $w_b$ & 同上 & 实现事实 \\
A 流带指示函数 & \code{activation}$\to$\code{gmm2} 的中间结果落点 & \textbf{编排选择} \\
\bottomrule
\end{tabular}
\end{center}

\noindent 中间结果落片上时 $\mathbb{1}[\text{物化}]=0$, 且 GMM2 必须与产它的 activation
\textbf{同核} (结果经固定管道直给配对的向量核, 该通路只在绑定对内存在)。因此这一项
同时改\textbf{公式}与\textbf{图结构}, 是``成本''与``约束''分开建模的典型例子。

\subsection{combine}
\begin{equation}
T_{\text{comb}} =
  \underbrace{\frac{m\,(e_{\text{in}} n + \mu)}{\BW_{\text{local}}}}_{\text{读回}}
+ \underbrace{\frac{(m-r)\, e_{\text{out}} n}{\BW_{\text{local}}}}_{\text{本卡写}}
+ \underbrace{\frac{r\, e_{\text{out}} n}{\BW_{\text{remote}}}}_{\text{跨卡写}}
+ \underbrace{m\,\alpha\left(\frac{s}{m}\right)^{\!\gamma}}_{\text{落点跨度}} ,
\label{eq:combine}
\end{equation}
$r$ 为本窗中目的卡不等于本卡的行数 (由路由精确给出), $\mu$ 为每行路由元数据字节,
$s$ 为落点跨度。$r$ 是必填参数: 跨卡写占实测的约 95\%, 缺了它公式低估一个数量级。
$\alpha$ 缺省为 0, 即不声称存在跨度效应 —— 实测只有两点且呈超线性 ($m$ 比 3.56 倍而
时长比 6.62 倍), 按字节与按每行固定开销都解释不了, 两点不足以定 $\gamma$。

\subsection{尾段}
六个事件串成一条链, 五项为常数, 一项按字节:
\begin{equation}
T_{\text{unpermute}} = \frac{\text{tokens}\,(\text{topk}\cdot h\cdot 2 + h\cdot 2)}{\BW_{\text{agg}}} .
\label{eq:unpermute}
\end{equation}
这条链属于某一份具体实现, 不属于框架。

%======================================================================
\section{图的约束}
\label{sec:graph}

一个事件是五元组
\begin{equation}
e = \big(R(e),\ (P(e),\lambda),\ (A(e),L(e)),\ c(e),\ d(e)\big),
\label{eq:event}
\end{equation}
语义如下。五个字段就是这套 IR 的表达力上限。

\begin{center}
\begin{tabular}{@{}llp{6.2cm}@{}}
\toprule
字段 & 符号 & 语义与时间跨度 \\
\midrule
独占资源 & $R(e)$ & 执行期独占; 可同时占多个; 写成池占位时在派发时刻才绑定 \\
前驱与边延迟 & $P(e), \lambda$ & 偏序 \\
计数额度 & $A(e), L(e)$ & \textbf{跨事件}持有: 在 $A$ 处取、在另一事件 $L$ 处还 \\
共位锚点 & $c(e)$ & 必须与某事件同核 \\
时长 & $d(e)$ & 第 \ref{sec:cost} 节的结果 \\
\bottomrule
\end{tabular}
\end{center}

\paragraph{为什么双缓冲必须用计数额度而非边.} 依赖边表达``$B$ 不早于 $A$ 结束''。双缓冲
要表达的是``第 $i+2$ 个生产者不能开始, 直到第 $i$ 个消费者归还缓冲槽'', 约束的两端随
在途数滑动, 不是固定的一对事件, 画不成静态边。

%======================================================================
\section{调度规则}
\label{sec:sched}

对每个就绪事件, 最早可行时刻
\begin{equation}
t_{\text{base}}(e) = \max\left\{
  \max_{p\in P(e)}\big[\mathrm{end}(p) + \lambda(e,p)\big],\
  \max_{r\in R(e)} \mathrm{free}(r) \right\},
\label{eq:tbase}
\end{equation}
再在额度上做准入搜索
\begin{equation}
\mathrm{start}(e) = \min\big\{\, t \ge t_{\text{base}}(e)\ :\ A(e)\ \text{的额度在}\ t\ \text{可满足} \,\big\}.
\label{eq:start}
\end{equation}
具体资源对 $R(e)$ 取 $\max$ (都必须空), 池占位取 $\min$ (哪个成员先空用哪个, 即晚绑定)。
\eqref{eq:start} 是在\textbf{归还时刻}上搜索而不是排队, 因此一个事件可以插到更早的
归还点, 不必等前面全部走完。就绪队列按 $(\mathrm{start}, \mathrm{order}, \mathrm{name})$
取最小者提交。模型的输出即该调度的 makespan
\begin{equation}
\widehat{T}(I) = C_{\max}\big(\sigma(I)\big) = \max_e \mathrm{end}(e),
\label{eq:makespan}
\end{equation}
其中 $\sigma$ 为本节的调度规则。\eqref{eq:instance} 与 \eqref{eq:makespan} 把两个独立的
建模决定分开: 实例怎么构造 (第 \ref{sec:shape}--\ref{sec:graph} 节), 以及调度规则是什么
(本节)。

\paragraph{归类与界.} 这是 $P|\mathrm{prec}|C_{\max}$ 的推广 (多资源同时独占、跨事件计数
额度、共位约束), NP-hard。表调度给出\textbf{可行解}而非最优解。

\paragraph{时序异常.} 表调度有已知的非单调性: 缩短某事件、放松某约束、增加核数, 都可能
使 makespan 变大。本模型中实测到过加入一条 $0.01\ \mu s$ 的边使 makespan 下降 2.81\% 的
例子。因此\textbf{贪心是模型机制而非硬件行为} (被建模的那份内核按滚动游标在编译期分块,
硬件不重分配工作), 且 $|\Delta\widehat{T}|$ 很小时其符号不可信。

%======================================================================
\section{物理下界}
\label{sec:bounds}

\subsection{算法事实}
下列两个量只由形状决定, 与 tile 切法、分核方式、编排\textbf{全部无关}:
\begin{align}
\MAC_{\text{total}} &= \MAC_{\text{gmm1}} + \MAC_{\text{gmm2}} = 3.058\times 10^{10},\\
\mathrm{bytes}_{\mathrm{GM}\to\mathrm{L1}} &= A_1 + B_1 + A_2 + B_2 = 216.5\ \mathrm{MB},
\end{align}
字节按``每个元素至少跨 GM$\to$L1 一次''计。

\subsection{三条界}
\begin{align}
\LB_{\text{compute}} &= \frac{\MAC_{\text{total}}}{R_{\text{cube}}\, P} = 40.45\ \mu s, \label{eq:lb-c}\\
\LB_{\text{bw}} &= \frac{\mathrm{bytes}}{\min(\BW_{\text{core}} P,\ \BW_{\text{agg}})} = 149.01\ \mu s, \label{eq:lb-b}\\
\LB_{\text{dep}} &= \sum_{s\in\text{chain}} \min_{e\in s} d(e) = 84.45\ \mu s, \label{eq:lb-d}\\
\LB &= \max(\LB_{\text{compute}}, \LB_{\text{bw}}, \LB_{\text{dep}}) = 149.01\ \mu s .
\end{align}
\eqref{eq:lb-c} 与 \eqref{eq:lb-b} 是\textbf{面积论证} (总量除以吞吐),
\eqref{eq:lb-d} 是\textbf{偏序论证} (必经链上各段的最小一份), 三者相互独立, 故必须都留。

\subsection{双向可证伪}
\begin{equation}
\widehat{T} < \LB \ \Longrightarrow\ \text{模型漏算了某项代价},
\end{equation}
\begin{equation}
T_{\text{实测}} < \LB \ \Longrightarrow\ \text{输入事实有错 (下界不含任何编排假设)} .
\end{equation}
本模型当前存在一个后者的未解实例: 实测最小值 $124.977\ \mu s$ 比 \eqref{eq:lb-b} 低
16.1\%。候选解释是权重量化位宽、板级规格档位, 或测量窗口口径不同。

%======================================================================
\section{空闲分解与守恒}
\label{sec:idle}

对某个角色池内的每个核,
\begin{equation}
\text{horizon} = \text{busy} + \text{forced} + \text{avoidable},
\qquad 200.91 + 51.42 + 49.60 = 301.93\ \mu s .
\label{eq:idle}
\end{equation}
\code{forced} 指该时刻没有就绪未开始的事件 (不可消除, 只能改图);
\code{avoidable} 指有就绪事件却仍有核空着, 在池化绑定下这是守恒性违规
\begin{equation}
\text{avoidable} > 0 \ \Longrightarrow\ \texttt{WorkConservationViolation}.
\end{equation}

\paragraph{汇报口径.} 各核空闲量\textbf{求和}得到的``核$\cdot\mu s$''是工作量单位而非时间,
核是并行的, 把它当总时长会误读。应报墙钟量: 本例时间轴上有核可避免空闲的窗口并集为
$86.02\ \mu s$ (3 个不相连窗口), 期间平均 16.1 个核 (共 28) 同时空闲。

\paragraph{avoidable 是上界而非账.} 有活就绪不等于搬过去就能缩短 makespan: 本例开启
晚绑定实际回收 $37.07\ \mu s$, 为上界的 43\%, 因为最大的那个窗口右端顶在关键路径节点上。

%======================================================================
\section{未建模的部分}
\label{sec:gaps}

\begin{center}
\begin{tabular}{@{}lp{8.2cm}@{}}
\toprule
项 & 状态 \\
\midrule
带宽争用 & 事件申报字节只进访存量统计, 不参与准入、不影响时长 \\
计数门控常数 & 明确未建模; 常数仅作复现记录。要建模应在尾段发一个真实事件 \\
引擎队列深度 & \textbf{刻意不设为参数}: 没有发射开销或在途计数的物理后果, 四类形状扫过任何取值结果逐位相同 \\
动态取活开销 & 缺省 0 且标注为假设 —— 0 不表示``没有代价'', 而是``本模型未声称代价是多少'' \\
落点跨度系数 & $\alpha=0$, 两个实测点不足以定 $\gamma$ \\
片上容量上界 & $\text{items}\times\text{tile} \le \mathrm{L1/UB}$ 未强制 \\
\bottomrule
\end{tabular}
\end{center}

\paragraph{参数的存在条件.} 一个参数必须有可建模的物理后果才配存在; 否则扫描会得到
``改它没用''的假结论。常数按出处分桶 (实现 / 实测 / 规格 / 算法 / 假设), 未标定输入分
两类: 有实测区间的可传播成结论的区间, 连范围都没有的\textbf{不许编一个范围} ——
那是把无知包装成精度。

%======================================================================
\section{适用的推断}
\label{sec:use}

\begin{center}
\begin{tabular}{@{}lll@{}}
\toprule
推断 & 可靠性 & 理由 \\
\midrule
同一声明下换一个旋钮的 $\Delta\widehat{T}$ & 可靠 & 两次运行共享全部未标定量, 系统误差同向抵消 \\
瓶颈归因 (哪条下界绑定) & 可靠 & 下界只依赖形状与硬件常数 \\
绝对时间 $\widehat{T}$ & 不可靠 & 直接暴露于每个假设类常数 \\
$|\Delta|$ 很小时的符号 & 不可靠 & 第 \ref{sec:sched} 节的时序异常 \\
\bottomrule
\end{tabular}
\end{center}

\noindent 因此本模型的产品是 $\Delta\widehat{T}$: 面向算子开发人员, 在同一份声明下改变
编排算法或 tiling 粒度并重跑, 差值即该选择的代价; 两次运行通过同一组护栏
(第 \ref{sec:bounds}--\ref{sec:idle} 节), 故差值可比而非模拟器的假象。

\end{document}
